\section{Segunda línea principal: la capacidad del modelo y la ingeniería de sistemas no son una relación de sustitución}

\subsection{De LM a LM system, y luego a Agent system}
Un punto muy valioso de la discusión es que los panelistas no opusieron el «bando del modelo» y el «bando del sistema». Al contrario, consideraron que el sistema de Agent actual es una extensión natural del LM system: además del modelo, se necesita infra, memory, workflow personalizado, interfaz de interacción y encapsulamiento de herramientas. En otras palabras, un modelo fuerte no hace desaparecer la ingeniería de sistemas, solo cambia su punto focal.

\textit{«No es solo un modelo el que presta el servicio, sino todo un sistema centrado en el modelo grande el que nos lo presta.»}

\begin{importantbox}{Por qué Cursor y Claude Code se volvieron ejemplos frecuentes}
El atractivo de estos productos no proviene solo de la capacidad de coding del modelo en sí, sino también de todo el sistema de interacción:
\begin{itemize}
  \item cómo se inyecta el contexto del editor;
  \item cómo se gestionan los permisos de archivos y comandos;
  \item cómo se retroalimentan al modelo los resultados de las llamadas a herramientas;
  \item cómo el usuario toma el control del proceso en los puntos clave.
\end{itemize}
Esto muestra que un «buen producto de Agent» es tanto un problema de modelo como un problema de diseño de sistema de producto.
\end{importantbox}

\subsection{Los límites entre memory, workflow y tool space se están moviendo}
En los prompting agent orientados a 2024, muchas capacidades dependían sobre todo de un workflow externo; para 2025, los panelistas observaron que cada vez más capacidades comenzaban a migrar hacia el interior del modelo, incluyendo hábitos de coding, estrategias de computer use, ciertos patrones de reasoning, e incluso parte de las capacidades de memory y tool use. Al mismo tiempo, el sistema externo no desapareció, sino que pasó a ocuparse de nuevos problemas: un control de permisos más complejo, una personalización de usuario más fina y una gestión de estado de más largo plazo.

\begin{table}[H]
\centering
\begin{tabular}{p{3.0cm}p{4.3cm}p{4.3cm}}
\toprule
\textbf{Capa del sistema} & \textbf{Práctica típica en 2024} & \textbf{Cambio observado en la discusión de 2025} \\
\midrule
Interior del modelo & Chat general, poco tool use & Comienza a interiorizar coding, GUI, multi-agent, reasoning \\
Memory & Recuperación externa o caché temporal & Se discute cómo interiorizarla parcialmente, pero el estado de largo plazo sigue dependiendo del sistema \\
Workflow & Principalmente orquestación de prompts & Transición gradual hacia una exploración autónoma y un aprendizaje de estrategias más capaces \\
Tool space & Pocas herramientas fijas & El espacio de acción se amplía, el entorno es más complejo y las restricciones son más determinantes \\
Interfaz de producto & Chat UI simple & Acoplamiento profundo con editores, dispositivos móviles y puntos de entrada de sistemas empresariales \\
\bottomrule
\end{tabular}
\caption{Un juicio implícito en el panel: el límite entre el modelo y el sistema se redibuja constantemente}
\end{table}

\subsection{Una vez reforzada la capacidad, el sistema asume el «último kilómetro»}
Los panelistas coincidieron en su juicio sobre la siguiente etapa: una vez que la capacidad agentic interna del modelo se refuerce hasta cierto punto, el foco de la investigación volverá a la capa externa del sistema. La razón no es complicada: un modelo fuerte solo puede elevar el límite superior, pero las aplicaciones empresariales necesitan estabilidad, aislamiento de permisos, personalización, seguridad e interfaces de despliegue. Incluso cuando el modelo es muy fuerte, estos problemas deben ser asumidos por el sistema.

\begin{warningbox}{No es realista interiorizar todas las capacidades en el modelo}
Si el environment, la memory, los permisos y las API externas dependen por completo de que el modelo los «invente» por su cuenta, el sistema perderá capacidad de control. Los panelistas no defendieron una interiorización total, sino que insistieron en que «donde haya un punto débil, ahí se debe reforzar primero». Se trata de un juicio marcadamente ingenieril, y no de una elección dogmática de architecture.
\end{warningbox}

\subsection{Resumen de la sección}
El desarrollo del Agent no consiste en que «el modelo sustituya al sistema», sino en que ambos se turnan como el punto débil. El foco de la investigación en 2025 se inclina más hacia inyectar capacidad dentro del modelo; pero en cuanto la capacidad se refuerza, memory, tooling, workflow, los permisos y la interfaz de conversión en producto vuelven a convertirse en el cuello de botella. Este vaivén es, en sí mismo, la trayectoria de maduración del sistema de Agent.

